\section{10. Limitations}

\subsection{10.1 Short-Sequence Silencing}

At $L = 128$ (our test length), FUSIONATTN uses the \textbf{short path} and 
provides no sparsity benefit. All our reported results reflect this — they 
are essentially LoRA-only results.

\subsection{10.2 Long Path Untested}

We did not test at $L > 128$ due to dataset and model constraints. Long-path 
behavior is extrapolated.

\subsection{10.3 Not a Compression Method}

Parameter count increases by ~1.17\%.

\subsection{10.4 Gate MLP Overhead}

The gate MLP adds ~0.20M parameters per model but is unused at 
L ≤ 128.

\subsection{10.5 Task Coverage}

Only SST-2 at $L = 128$. Generation tasks, QA, and long-context tasks 
untested.

\subsection{10.6 Model Coverage}

Three models (66.9M–124.6M).

\subsection{10.7 Hardware}

Only RTX 4070 Laptop.

\subsection{10.8 Causal Decoder Block Attention}

Block attention for decoders requires within-block causal masking. This is 
implemented but not extensively tested.

\subsection{10.9 No Formal Bound}

No bound on approximation error of top-K block selection.

\subsection{10.10 Single Seed}

Standard error ~1.69\%.

\subsection{10.11 Critical Bug: `raise RuntimeError` in inject}

The \texttt{inject_fusion_attn} function has:

\begin{lstlisting}
except Exception:
    raise RuntimeError("copy_weights failed")
\end{lstlisting}

This kills injection on any failure. Recommend changing to \texttt{return None}.

\subsection{10.12 Code Quality}

\begin{itemize}
\item Research-grade
\end{itemize}
\begin{itemize}
\item No unit tests for edge cases
\end{itemize}
\begin{itemize}
\item Broad exception handling
\end{itemize}

\subsection{10.13 Ethical Considerations}

Dual-use. See Appendix E.